\documentclass[10pt,a4paper]{amsart}
\usepackage[margin=19mm]{geometry}
\usepackage{amsmath,amssymb,mathtools}
\usepackage[scr=rsfs,cal=euler]{mathalfa}
\usepackage{xfrac}
\usepackage[unicode,hidelinks,bookmarks=false]{hyperref}
\newcommand{\ie}{i.e.\ }
\newcommand{\cf}{cf.\ }
\newcommand{\resp}{respectively\ }
\newcommand{\Subsection}{\subsection}
\providecommand{\nobreakdash}{\nobreak\hbox{-}}
\setlength{\emergencystretch}{2em}
\multlinegap0pt
\usepackage{amssymb}
\usepackage{mathtools}
\newtheorem{assumption}{Assumption}
\newtheorem{prop}{}
\def\R{{\mathbf R}}% real numbers
\def\T{{\mathbf T}}% torus
\def\d{{\partial}}
\DeclareMathOperator{\diver}{div}
\title{Initial layer analysis of relaxation-time limit of the collisional QHD}
\author{Paolo Antonelli, Pierangelo Marcati, Hao Zheng}
\date{Source: arXiv:2607.03773v1 (CC BY 4.0)}
\begin{document}
\maketitle

\noindent\emph{Source and licence.} Initial layer analysis of relaxation-time limit of the collisional QHD. Version arXiv:2607.03773v1 (CC BY 4.0), distributed by arXiv under the Creative Commons Attribution 4.0 International licence, http://creativecommons.org/licenses/by/4.0/, as recorded in the arXiv licence metadata for this exact version and shown on the abstract page https://arxiv.org/abs/2607.03773v1. Full licence notice: \texttt{licenses/arxiv-2607.03773-CC-BY-4.0.txt}.\par\smallskip

\noindent\textit{Editorial scope.} Counted unit: the article's prop eq:dtEr (source0.tex lines 777--788). The article's complete original proof of it is delivered with it (source0.tex lines 790--908, one \texttt{proof} environment of 119 lines). No mathematics is written, repaired or re-derived: the statement and the proof are the article's own lines, in the article's own notation, and no retained line is substituted or re-flowed.  Retained uncounted as the source-local context the proof needs: lines 127--133; lines 319--327; lines 585--592; lines 676--678; lines 680--685; lines 687--689; lines 691--696; lines 713--715; lines 744--763.  Nothing was omitted from any retained span, so the package carries no omission record.  No retained line contains a citation, so the wrapper carries no bibliography and no bibliography entry.  No retained line contains a figure environment, an image-inclusion command or drawing code, so no image is delivered and the package declares no asset dependency.  Editorial formatting only, in addition: an AMS \texttt{amsart} preamble that restates the article's own declarations and macros used by the retained lines, namely the environments \texttt{assumption}, \texttt{prop} on separate counters, together with the standard packages those lines need.  The retained environments print in this document as Assumption 1, Proposition 1 in order of appearance, whereas the article prints the same items with the numbering of its own full document.  The complete native source of the article as distributed in the arXiv e-print for this exact version is bundled unchanged as \texttt{sources/204418/source0.tex}.
\begin{equation}\label{eq:QHD_rs_intro}
\left\{\begin{aligned}
&\d_{t}\rho_\tau+\diver J_\tau=0\\
&\tau^2\d_{t} J_\tau+\tau^2\diver\left(\frac{J_\tau\otimes J_\tau}{\rho_\tau}\right)+\nabla p(\rho_\tau)+\rho_\tau\nabla V_\tau=\frac{1}{2}\rho_\tau\nabla\left(\frac{\triangle\sqrt\rho_\tau}{\sqrt\rho_\tau}\right)-J_\tau\\
&-\triangle V_\tau=\rho_\tau-\mathcal{C}(x),\quad (\rho_\tau,J_\tau)(0,x)=(\rho_0,\tau^{-1}J_0)(x).
\end{aligned}\right.
\end{equation}

\begin{equation}\label{eq:opB}
\begin{aligned}
\mathcal{B}(\rho)=&\frac14\nabla\triangle r-\frac14\diver\left[\frac{2(\nabla\rho_0\otimes \nabla r)^s+\nabla \eta_\tau\otimes \nabla\eta_\tau}{\rho}\right]+\frac14\diver\left[\frac{r\nabla \rho_0\otimes\nabla\rho_0}{\rho_0^2}\right]\\
&-\frac{1}{2\tau^2}\diver\left[(\nabla \rho_0\otimes\nabla\rho_0) \int_0^{\tau \eta_\tau}\frac{\tau \eta_\tau-g}{(\rho_0+g)^3}dg\right]+\frac{1}{2}\diver\left[\frac{\eta_\tau(\nabla \rho_0\otimes\nabla\rho_1)^s}{\rho_0^2}\right]\\
&-\frac{1}{\tau}\diver\left[(\nabla\rho_0\otimes\nabla\rho_1)^s\int_0^{\tau \eta_\tau}\frac{\tau \eta_\tau-g}{(\rho_0+g)^3}dg\right]\\
&-\nabla[p'(\rho_0)r]-\tau^{-2}\nabla\int_0^{\tau \eta_\tau} p''(\rho_0+g)(\tau \eta_\tau-g)dg\\
&-r\nabla V(\rho_0)-\rho_0\nabla V(r)-\eta_\tau\nabla V(\eta_\tau).
\end{aligned}
\end{equation}

\begin{equation}\label{eq:J_aprx}
\begin{aligned}
\tilde J_\tau(s,t)=&\tau^{-1}e^{-s}J_0-e^{-s}\left[\bar J(\rho_0)+(1-e^{-s})\diver\left(\frac{J_0\otimes J_0}{\rho_0}\right)\right]+\bar J(t)\\
&+2\tau e^{-s}(1-e^{-s}-s)\diver\left[\frac{J_0\otimes \bar J(\rho_0)}{\rho_0}\right]^s\\
&+2\tau e^{-s}(1+e^{-2s}-2e^{-s})\diver\left\{\left[\frac{J_0}{\rho_0}\otimes\diver\left(\frac{J_0\otimes J_0}{\rho_0}\right)\right]^s-\frac{J_0\otimes J_0\diver J_0}{2\rho_0^2}\right\}\\
&-\tau e^{-s}(s+1)\mathcal{A}(\rho_0,\diver J_0)+\tau J_{\mathcal{O},1}(t)
\end{aligned}
\end{equation}

\begin{equation}\label{eq:dtr2}
\d_t r-\diver \mathcal{S}_\rho+\diver\overline{\mathcal{R}}=0,
\end{equation}

\begin{equation}\label{eq:dtR2}
\begin{aligned}
\tau^2\d_t\overline{\mathcal{R}}+\tau^2\diver\left(\frac{J_\tau\otimes\overline{\mathcal{R}}}{\rho_\tau}\right)+&\tau^2\diver\left[\frac{\overline{\mathcal{R}}\otimes(\tau^{-1}e^{-s}J_0+\xi_\tau)}{\rho_\tau}\right]\\
=&\mathcal{S}_J+\mathcal{B}(\bar\rho,e^{-s}\diver J_0+\rho_{\mathcal{O},1},r)-\overline{\mathcal{R}},
\end{aligned}
\end{equation}

\begin{equation}\label{eq:S_rho}
\mathcal{S}_\rho=\tau^{-2}e^{-s}\left[\bar J(\rho_0)+(1-e^{-s})\diver\left(\frac{J_0\otimes J_0}{\rho_0}\right)\right].
\end{equation}

\begin{equation}\label{eq:S2}
\begin{aligned}
\mathcal{S}_J=&\tau^{-1} e^{-s}\mathcal{A}(\bar\rho,\diver J_0)+\tau^{-2}e^{-2s}\diver\left[(J_0\otimes J_0)(\rho_0^{-1}-\rho_\tau^{-1})\right]\\
&-\diver\left[\frac{2\tau^{-1}e^{-s}\xi_\tau\otimes J_0+\xi_\tau\otimes \xi_\tau}{\rho_\tau}\right]-\sum_{k=0}^1 \tau^{k}\d_t J_{\mathcal{O},k}.
\end{aligned}
\end{equation}

\begin{equation}\label{eq:xi}
\xi_\tau=-e^{-s}\left[\bar J(\rho_0)+(1-e^{-s})\diver\left(\frac{J_0\otimes J_0}{\rho_0}\right)\right]+\bar J(t)+\tau J_{\mathcal{O},1}.
\end{equation}

\begin{assumption}[Regularity assumption]\label{asmp:RA}
Assume that on the time interval $[0,T_0]$, the following conditions hold.
\begin{itemize}
\item[(1)] Positivity of densities: there exists a $\delta>0$ such that
\[
\inf_{t,x}\,\{\rho_0,\rho_\tau,\bar\rho\}\geq\delta >0.
\]

\item[(2)] Irrotationality: the initial velocity $v_0=\frac{J_0}{\rho_0}$ satisfies
\[
(\nabla v_0)^a=\frac12\left(\nabla v_0-\nabla v_0^T\right)=0.
\]

\item[(3)] Uniform Sobolev regularity: there exists a $M>0$ such that
\begin{align*}
\max&\left\{\|\rho_0\|_{H^{m+1}_x},\|J_0\|_{H^{m}_x},\|\rho_\tau\|_{L^\infty_x H^{m+1}_x},\right.\\
&\hspace{0.3cm}\left.\|J_\tau\|_{L^\infty_x H^{m}_x},\|\bar\rho\|_{L^\infty_x H^{m+1}_x},\|\rho_{\mathcal{O},1}\|_{L^\infty_x H^{m-1}_x}\right\}\leq M.
\end{align*}
\end{itemize}
\end{assumption}

\begin{prop}
Provided regularity Assumption \ref{asmp:RA}, the energy functional $E_\tau(r,\overline{\mathcal{R}})$ can be differentiable in time along the trajectory, and we have 
\begin{equation}\label{eq:dtEr}
\begin{aligned}
\frac{d}{dt}E_\tau(r,\overline{\mathcal{R}})+\int_{{\T^d}} \frac{|\overline{\mathcal{R}}|^2}{\rho_\tau}dx=&\int_{{\T^d}} \frac{\overline{\mathcal{R}}}{\rho_\tau}\cdot\mathcal{S}_J(t,x)dx-\frac{\tau^2}{2}\int_{\T^d}\frac{|\overline{\mathcal{R}}|^2}{\rho_\tau^2}\diver(\tau^{-1}e^{-s}J_0+\xi_\tau)dx\\
&+\int_{\T^d}\frac{\overline{\mathcal{R}}}{\rho_\tau}\cdot\left[-\tau^{-1} e^{-s}\left(v_0\otimes\nabla\left(\frac{\rho_0}{\rho_\tau}\right)\right)^a+\mathcal{S}_{rot}\right]\cdot(\tau^{-1}e^{-s}J_0+\xi_\tau)dx\\
&-\int_{\T^d}\frac{|\nabla r|^2}{2\rho_\tau^2}\d_t\rho_\tau dx+\int_{\T^d} \frac{\nabla\diver \mathcal{S}_\rho}{\rho_\tau}\cdot\nabla rdx\\
&+\int_{\T^d}\frac{\overline{\mathcal{R}}}{\rho_\tau}\cdot[\mathcal{B}_2(r,\nabla r)+\mathcal{B}_3(\nabla r)]dx,
\end{aligned}
\end{equation}
where $\mathcal{S}_{rot}$, $\mathcal{S}_J$, $\xi_\tau$, $\mathcal{S}_\rho$, $\mathcal{B}_2(r,\nabla r)$ and $\mathcal{B}_3(\nabla r)$ are given by \eqref{eq:S_rho}, \eqref{eq:S2}, \eqref{eq:xi}, \eqref{eq:S_curl}, \eqref{eq:B2} and \eqref{eq:B3} respectively.
\end{prop}

\begin{proof}

We start with equation \eqref{eq:dtR2} by multiplying \eqref{eq:dtR2} by $\overline{\mathcal{R}}/\rho_\tau$. For the advection term, we notice that
\begin{align*}
\tau^2\int_{\T^d} \frac{\overline{\mathcal{R}}}{\rho_\tau}\cdot\diver\left(\frac{J_\tau\otimes\overline{\mathcal{R}}}{\rho_\tau}\right)dx=&\tau^2\int_{\T^d} \frac{|\overline{\mathcal{R}}|^2}{\rho_\tau^2}\diver J_\tau dx+\frac{\tau^2}{2}\int_{\T^d} J_\tau\cdot\nabla\left(\frac{|\overline{\mathcal{R}}|^2}{\rho_\tau^2}\right)dx\\
=&-\frac{\tau^2}{2}\int_{\T^d} \frac{|\overline{\mathcal{R}}|^2}{\rho_\tau^2}\d_t\rho_\tau dx
\end{align*}
then we have
\begin{align*}
&\frac{\tau^2}{2}\frac{d}{dt}\int_{\T^d} \frac{|\overline{\mathcal{R}}|^2}{\rho_\tau}dx+\int_{\T^d} \frac{|\overline{\mathcal{R}}|^2}{\rho_\tau}dx\\
=&-\tau^2\int_{\T^d} \frac{\overline{\mathcal{R}}}{\rho_\tau}\cdot\diver\left[\frac{\overline{\mathcal{R}}\otimes(\tau^{-1}e^{-s}J_0+\xi_\tau)}{\rho_\tau}\right]dx+\int_{\T^d} \frac{\overline{\mathcal{R}}}{\rho_\tau}\cdot\mathcal{S}_J(t,x)dx\\
&+\int_{\T^d} \frac{\overline{\mathcal{R}}}{\rho_\tau}\cdot\mathcal{B}(\bar\rho,e^{-s}\diver J_0+\rho_{\mathcal{O},1},r)dx.
\end{align*}


The integral
\[
-\tau^2\int_{\T^d} \frac{\overline{\mathcal{R}}}{\rho_\tau}\cdot\diver\left[\frac{\overline{\mathcal{R}}\otimes(\tau^{-1}e^{-s}J_0+\xi_\tau)}{\rho_\tau}\right]dx
\]
can be computed by integrating by parts as
\begin{align*}
=&\tau^2\int_{\T^d} \d_j\left(\frac{\overline{\mathcal{R}}_k}{\rho_\tau}\right)\frac{\overline{\mathcal{R}}_j\otimes(\tau^{-1}e^{-s}J_0+\xi_\tau)_k}{\rho_\tau}dx\\
=&\tau^2\int_{\T^d} \d_k\left(\frac{\overline{\mathcal{R}}_j}{\rho_\tau}\right)\frac{\overline{\mathcal{R}}_j(\tau^{-1}e^{-s}J_0+\xi_\tau)_k}{\rho_\tau}dx\\
&+\tau^2\int_{\T^d} \left[\d_j\left(\frac{\overline{\mathcal{R}}_k}{\rho_\tau}\right)-\d_k\left(\frac{\overline{\mathcal{R}}_j}{\rho_\tau}\right)\right]\frac{\overline{\mathcal{R}}_j(\tau^{-1}e^{-s}J_0+\xi_\tau)_k}{\rho_\tau}dx\\
=&\frac{\tau^2}{2}\int_{\T^d} \nabla\left(\frac{|\overline{\mathcal{R}}|^2}{\rho_\tau^2}\right)\cdot(\tau^{-1}e^{-s}J_0+\xi_\tau) dx+\tau^2\int_{\T^d} \left(\nabla\frac{\overline{\mathcal{R}}}{\rho_\tau}\right)^a:\frac{\overline{\mathcal{R}}\otimes(\tau^{-1}e^{-s}J_0+\xi_\tau)}{\rho_\tau}dx\\
=&-\frac{\tau^2}{2}\int_{\T^d}\frac{|\overline{\mathcal{R}}|^2}{\rho_\tau^2}\diver(\tau^{-1}e^{-s}J_0+\xi_\tau)dx+\tau^2\int_{\T^d} \left(\nabla\frac{\overline{\mathcal{R}}}{\rho_\tau}\right)^a:\frac{\overline{\mathcal{R}}\otimes(\tau^{-1}e^{-s}J_0+\xi_\tau)}{\rho_\tau}dx,
\end{align*}
where $A^a=\frac12(A-A^T)$ is the anti-symmetric part of matrix $A\in\R^{d\times d}$. 
To treat the integral of the curl-matrix, we recall that the irrotaitonality of the velocity is preserved by the trajectory of \eqref{eq:QHD_rs_intro}, namely 
\[
(\nabla v_\tau)^a=\left(\nabla\frac{J_\tau}{\rho_\tau}\right)^a=0,
\]
by definition of $\overline{\mathcal{R}}$ and \eqref{eq:J_aprx}, it follows that
\begin{align*}
0=(\nabla v_\tau)^a=&\tau^{-1} e^{-s}\left(\nabla\frac{J_0}{\rho_\tau}\right)^a-e^{-s}\left[\nabla\frac{\bar J(\rho_0)}{\rho_\tau}+\frac{(1-e^{-s})}{\rho_\tau}\nabla\diver\left(\frac{J_0\otimes J_0}{\rho_0}\right)\right]^a\\
&+\left(\nabla\frac{\bar J(t)}{\rho_\tau}\right)^a+\tau \left(\nabla\frac{J_{\mathcal{O},1}}{\rho_\tau}\right)^a+\tau^2 \left(\nabla\frac{\overline{\mathcal{R}}}{\rho_\tau}\right)^a.
\end{align*}
Therefore
\begin{align*}
\tau^2 \left(\nabla\frac{\overline{\mathcal{R}}}{\rho_\tau}\right)^a
=-\tau^{-1} e^{-s}\left[\nabla\left(v_0\frac{\rho_0}{\rho_\tau}\right)\right]^a+\mathcal{S}_{rot}(t,x),
\end{align*}
where 
\begin{equation}\label{eq:S_curl}
\begin{aligned}
\mathcal{S}_{rot}(t,x)=&e^{-s}\left[\nabla\frac{\bar J(\rho_0)}{\rho_\tau}+\frac{(1-e^{-s})}{\rho_\tau}\nabla\diver\left(\frac{J_0\otimes J_0}{\rho_0}\right)\right]^a\\
&-\left(\nabla\frac{\bar J(t)}{\rho_\tau}\right)-\tau \left(\nabla\frac{J_{\mathcal{O},1}}{\rho_\tau}\right)
\end{aligned}
\end{equation}
and $\mathcal{S}_{rot}$ is uniformly bounded in $L^\infty_t$. For the first term in the right hand side of $\tau^2 \left(\nabla\frac{\overline{\mathcal{R}}}{\rho_\tau}\right)^a$, we further notice that by the irrotationality of $v_0$, we have
\[
\left[\nabla\left(v_0\frac{\rho_0}{\rho_\tau}\right)\right]=\left(v_0\otimes\nabla\left(\frac{\rho_0}{\rho_\tau}\right)\right)^a,
\]
thus we obtain that
\begin{align*}
&\tau^2\int_{\T^d} \left(\nabla\frac{\overline{\mathcal{R}}}{\rho_\tau}\right)^a:\frac{\overline{\mathcal{R}}\otimes(\tau^{-1}e^{-s}J_0+\xi_\tau)}{\rho_\tau}dx\\
=&\int_{\T^d}\frac{\overline{\mathcal{R}}}{\rho_\tau}\cdot\left[-\tau^{-1} e^{-s}\left(v_0\otimes\nabla\left(\frac{\rho_0}{\rho_\tau}\right)\right)^a+\mathcal{S}_{rot}\right]\cdot(\tau^{-1}e^{-s}J_0+\xi_\tau)dx.
\end{align*}


Now we consider the integral containing the operator $\mathcal{B}$. Recall from \eqref{eq:opB} that we have
\begin{align*}
&\mathcal{B}(\bar\rho,e^{-s}\diver J_0+\rho_{\mathcal{O},1},r)\\
=&\frac14\nabla\triangle r-\frac14\diver\left[\frac{2(\nabla\bar\rho \otimes\nabla r)^s+\nabla \eta_\tau\otimes \nabla \eta_\tau}{\rho_\tau}\right]+\mathcal{B}_2(r,\nabla r).
\end{align*}
where
\begin{equation}\label{eq:B2}
\begin{aligned}
\mathcal{B}_2(r,\nabla r)=&\frac14\diver\left[\frac{r\nabla \bar\rho\otimes\nabla \bar\rho }{\bar\rho^2}\right]-\frac{1}{2\tau^2}\diver\left[\nabla \bar\rho\otimes\nabla \bar\rho\int_0^{\tau \eta_\tau}\frac{\tau \eta_\tau-g}{(\bar\rho+g)^3}dg\right]\\
&+\frac12\diver\left[\frac{(\nabla \bar\rho\otimes\nabla(e^{-s}\diver J_0+\rho_{\mathcal{O},1}))^sr_\tau}{\bar\rho^2}\right]\\
&-\frac{1}{\tau}\diver\left[(\nabla \bar\rho\otimes\nabla(e^{-s}\diver J_0+\rho_{\mathcal{O},1}))^s\int_0^{\tau \eta_\tau}\frac{\tau \eta_\tau-g}{(\bar\rho+g)^3}dg\right]\\
&-\nabla [p'(\bar\rho)r]-\tau^{-2}\nabla\int_0^{\tau \eta_\tau} p''(\bar\rho+g)(\tau \eta_\tau-g)dg\\
&-r\nabla V(\bar\rho)-\bar\rho\nabla V(r)-\eta_\tau\nabla V(\eta_\tau)
\end{aligned}
\end{equation}
and
\[
\eta_\tau=e^{-s}\diver J_0+\rho_{\mathcal{O},1}+\tau r.
\]
We first treat the first two leading order terms in the right hand side of $\mathcal{B}$,
\begin{align*}
\int_{\T^d} \frac{\overline{\mathcal{R}}}{\rho_\tau}\cdot\nabla\triangle rdx-\int_{\T^d} \frac{\overline{\mathcal{R}}}{\rho_\tau}\cdot\diver\left[\frac{2(\nabla\bar\rho \otimes\nabla r)^s+\nabla \eta_\tau\otimes \nabla \eta_\tau}{\rho_\tau}\right]dx.
\end{align*}
By integrating by parts and \eqref{eq:dtr2}, we have
\begin{align*}
\int_{\T^d} \frac{\overline{\mathcal{R}}}{\rho_\tau}\cdot\nabla\triangle rdx=&-\int_{\T^d} \frac{\diver \overline{\mathcal{R}}}{\rho_\tau}\triangle rdx+\int_{\T^d} \frac{\overline{\mathcal{R}}}{\rho_\tau^2}\cdot\nabla\rho_\tau\triangle r dx\\
=&\int_{\T^d} \frac{\nabla\diver \overline{\mathcal{R}}}{\rho_\tau}\cdot\nabla rdx-\int_{\T^d} \frac{\diver \overline{\mathcal{R}}}{\rho_\tau^2}\nabla r\cdot\nabla\rho_\tau dx+\int_{\T^d} \frac{\overline{\mathcal{R}}}{\rho_\tau^2}\cdot\nabla\rho_\tau\triangle r  dx\\
=&-\frac{d}{dt}\int_{\T^d}\frac{|\nabla r|^2}{2\rho_\tau}dx-\int_{\T^d}\frac{|\nabla r|^2}{2\rho_\tau^2}\d_t\rho_\tau dx+\int_{\T^d} \frac{\nabla\diver \mathcal{S}_\rho}{\rho_\tau}\cdot\nabla rdx\\
&-2\int_{\T^d} \frac{\overline{\mathcal{R}}\cdot\nabla\rho_\tau}{\rho_\tau^3}\nabla r\cdot\nabla\rho_\tau dx+\int_{\T^d} \frac{\overline{\mathcal{R}}}{\rho_\tau^2}\cdot(I_d\triangle+\nabla^2)r\cdot\nabla\rho_\tau dx\\
&+\int_{\T^d} \frac{\overline{\mathcal{R}}}{\rho_\tau^2}\cdot\nabla^2\rho_\tau\cdot\nabla r dx,
\end{align*}
where $I_d$ is the $d-$dimensional identity matrix. By the expansion of $\rho_\tau$, we notice that
\begin{align*}
&-\int_{\T^d} \frac{\overline{\mathcal{R}}}{\rho_\tau}\cdot\diver\left[\frac{2(\nabla\bar\rho\otimes \nabla r)^s+\nabla \eta_\tau\otimes \nabla \eta_\tau}{\rho_\tau}\right]dx\\
=&\int_{\T^d} \nabla\rho_\tau\cdot\left[\frac{2(\nabla\bar\rho\otimes \nabla r)^s+\nabla \eta_\tau\otimes \nabla \eta_\tau}{\rho_\tau^2}\right]\cdot\frac{\overline{\mathcal{R}}}{\rho_\tau} dx\\
&-\int_{\T^d} \frac{\overline{\mathcal{R}}}{\rho_\tau}\cdot\frac{\triangle\bar\rho\nabla r+\triangle (e^{-s}\diver J_0+\rho_{\mathcal{O},1}+\tau r)\nabla\eta_\tau}{\rho_\tau}dx\\
&-\int_{\T^d} \frac{\overline{\mathcal{R}}}{\rho_\tau}\cdot\frac{\nabla\bar\rho\cdot\nabla^2 r+\nabla\eta_\tau\cdot\nabla^2(e^{-s}\diver J_0+\rho_{\mathcal{O},1}+\tau r)}{\rho_\tau}dx\\
&-\int_{\T^d} \frac{\overline{\mathcal{R}}}{\rho_\tau}\cdot\frac{\triangle r\nabla\bar\rho+\nabla r\cdot\nabla^2\bar\rho}{\rho_\tau}dx.
\end{align*}
Recall that $\rho_\tau=\bar\rho+\tau\eta_\tau$, then we can re-organize the integrands of the right hand side above as
\begin{align*}
=&\int_{\T^d} \nabla\rho_\tau\cdot\left[\frac{2(\nabla\bar\rho\otimes \nabla r)^s+\nabla \eta_\tau\otimes \nabla \eta_\tau}{\rho_\tau^2}\right]\cdot\frac{\overline{\mathcal{R}}}{\rho_\tau} dx\\
&-\int_{\T^d} \frac{\overline{\mathcal{R}}}{\rho_\tau}\cdot\frac{\nabla r\cdot(I_d\triangle+\nabla^2 )\bar\rho}{\rho_\tau}dx-\int_{\T^d} \frac{\overline{\mathcal{R}}}{\rho_\tau}\cdot\frac{\nabla \eta_\tau\cdot(I_d\triangle+\nabla^2)(e^{-s}\diver J_0+\rho_{\mathcal{O},1})}{\rho_\tau}dx\\
&-\int_{\T^d} \frac{\overline{\mathcal{R}}}{\rho_\tau}\cdot\frac{\nabla\rho_\tau\cdot(I_d\triangle+\nabla^2)r}{\rho_\tau}dx
\end{align*}
Summarizing the identities above and noticing that the integrals containing $(I_d\triangle+\nabla^2)r$ cancel, we obtain
\begin{align*}
&\int_{\T^d} \frac{\overline{\mathcal{R}}}{\rho_\tau}\cdot\nabla\triangle rdx-\int_{\T^d} \frac{\overline{\mathcal{R}}}{\rho_\tau}\cdot\diver\left[\frac{2(\nabla\bar\rho\otimes \nabla r)^s+\nabla \eta_\tau\otimes \nabla \eta_\tau}{\rho_\tau}\right]dx\\
=&-\frac{d}{dt}\int_{\T^d}\frac{|\nabla r|^2}{2\rho_\tau}dx-\int_{\T^d}\frac{|\nabla r|^2}{2\rho_\tau^2}\d_t\rho_\tau dx+\int_{\T^d} \frac{\nabla\diver \mathcal{S}_\rho}{\rho_\tau}\cdot\nabla rdx+\int_{\T^d}\frac{\overline{\mathcal{R}}}{\rho_\tau}\cdot\mathcal{B}_3(\nabla r)dx,
\end{align*}
where
\begin{equation}\label{eq:B3}
\begin{aligned}
\mathcal{B}_3(\nabla r)=&(\nabla\rho_\tau\cdot\nabla r)\nabla \rho_\tau+\nabla^2\rho_\tau\cdot\nabla r+(I_d\triangle+\nabla^2)\bar\rho\cdot\nabla r\\
&-\frac{\nabla \eta_\tau}{\rho_\tau}\cdot(I_d\triangle+\nabla^2)(e^{-s}\diver J_0+\rho_{\mathcal{O},1}),
\end{aligned}
\end{equation}
Summarizing the identities above, we obtain \eqref{eq:dtEr}.
\end{proof}

\end{document}
